% Generated draft: requires mathematical and source review.
\documentclass[11pt]{article}
\usepackage[margin=25mm]{geometry}
\usepackage{fontspec}
\setmainfont{Latin Modern Roman}
\usepackage{amsmath,amssymb,mathtools,unicode-math}
\setmathfont{Latin Modern Math}
\usepackage{enumitem,needspace,xurl,hyperref}
\setlength{\parindent}{0pt}
\setlength{\parskip}{5pt}
\setlength{\emergencystretch}{3em}
\newcommand{\N}{\mathbb{N}}
\newcommand{\Z}{\mathbb{Z}}
\newcommand{\Q}{\mathbb{Q}}
\newcommand{\R}{\mathbb{R}}
\newcommand{\C}{\mathbb{C}}
\newcommand{\F}{\mathbb{F}}
\newcommand{\A}{\mathbb{A}}
\newcommand{\PP}{\mathbb{P}}
\newcommand{\E}{\mathbb{E}}
\newcommand{\eps}{\varepsilon}
\newcommand{\dd}{\,\mathrm{d}}
\begin{document}
\section*{UnsolvedMath ID: 9400099}
\textbf{Source problem code:} catalogue entry
\subsection{Definitions and mathematical statement}
In the Kontsevich--Zagier integral convention, an ordinary period is a complex number whose real and imaginary parts are absolutely convergent integrals of rational functions with rational coefficients over semialgebraic subsets of some \(\R^d\) defined over \(\Q\). Let \(\mathcal P\) be this class. Exponential periods allow suitable algebraic/rational exponential integrands, schematically
\[
\int_D R(x)e^{Q(x)}\,\mathrm dx,
\]
with algebraic data and the required convergence conditions. The exact exponential-period category, including complex coefficients and growth conditions, must be fixed; write it \(\mathcal P_{\exp}\). The ordinary class is included by taking \(Q=0\).

The question asks for a description or membership criteria for
\[
\mathcal P\setminus\overline\Q
\qquad\text{and/or}\qquad
\mathcal P_{\exp}\setminus\overline\Q,
\]
where \(\overline\Q\) denotes the algebraic numbers. These are different classification questions, and the parenthetical exponential does not make their known examples interchangeable. It does not specifically request an algorithm or a classification restricted to a stated family of constants.

For example,
\[
\pi=\int_{x^2+y^2\le1}1\,\mathrm dx\,\mathrm dy
\]
is an ordinary period, while
\[
e=\int_{-1}^{\infty}e^{-x}\,\mathrm dx
\]
is an exponential period in the usual integral sense. Whether \(e\) is an ordinary period is not settled by that representation. Thus the research summary's tentative period discussion must distinguish ordinary from exponential membership, and must not suggest that \(\pi\)'s ordinary-period status is unknown.

The supplied summary invokes the Kontsevich--Zagier conjecture but gives no precise statement. Its phrase algebraicity of periods when equal is not a usable formulation: equality of period representations does not make all periods algebraic. No relation conjecture is adopted as a necessary-and-sufficient classification of transcendental constants. The source supplies a broad research direction, not a single explicit universal membership conjecture.
\subsection{Short English statement}
Describe which transcendental numbers are ordinary periods and which are exponential periods. Keep these classes separate: \(\pi\) is an ordinary period and \(e\) is an exponential period, while the supplied broad classification and conjectural discussion remain unresolved.
\subsection{Sources}
\begin{itemize}
\item Local UnsolvedMath registry, permanent ID \(9400099\), exact source code \texttt{catalogue entry}, ``Which transcendental numbers are (exponential) periods.''
\item Wikipedia, \emph{List of unsolved problems in mathematics}, pinned revision 1366636767, Diophantine approximation and transcendental number theory bullet 99, wikitext line 125; \url{https://en.wikipedia.org/w/index.php?title=List_of_unsolved_problems_in_mathematics&oldid=1366636767#Number_theory}. MediaWiki-source extraction; recorded access 2026-07-29.
\item Supplied OPEN-TRIAGE summary points to Wikipedia's ``Period (algebraic geometry)'' article and mentions Kontsevich--Zagier, without a complete primary conjecture citation or a specified exponential-period definition.
\end{itemize}
\end{document}
